\documentclass[10pt,a4paper]{amsart}
\usepackage[margin=19mm]{geometry}
\usepackage{amsmath,amssymb,mathtools}
\usepackage[scr=rsfs,cal=euler]{mathalfa}
\usepackage{xfrac}
\usepackage[unicode,hidelinks,bookmarks=false]{hyperref}
\newcommand{\ie}{i.e.\ }
\newcommand{\cf}{cf.\ }
\newcommand{\resp}{respectively\ }
\newcommand{\Subsection}{\subsection}
\providecommand{\nobreakdash}{\nobreak\hbox{-}}
\setlength{\emergencystretch}{2em}
\multlinegap0pt
\usepackage{bm}
\usepackage{bbm}
\usepackage{dsfont}
\usepackage{xspace}
\usepackage{cancel}
\usepackage{mathtools}
\usepackage{amsthm}
\makeatletter
\@ifundefined{thm}{\newtheorem{thm}{Thm}}{}
\@ifundefined{theorem}{\newtheorem{theorem}[thm]{Theorem}}{}
\makeatother
\title[Poisson structures on weak Sobolev loop spaces and applicati]{Poisson structures on weak Sobolev loop spaces and applications to integrable systems}
\author{Jean-Pierre Magnot}
\date{Source: arXiv:2510.19830v1 (CC BY 4.0)}
\begin{document}
\maketitle

\noindent\emph{Source and licence.} Poisson structures on weak Sobolev loop spaces and applications to integrable systems. Version arXiv:2510.19830v1 (CC BY 4.0), distributed under the Creative Commons Attribution 4.0 International licence, as recorded in the arXiv licence metadata for this exact version and shown on the abstract page https://arxiv.org/abs/2510.19830v1. Full rights notice: licenses/arxiv-2510.19830-CC-BY-4.0.txt.\par\smallskip

\noindent\textit{Editorial scope.} Counted unit: the Theorem whose source label is \texttt{th:presymp-Hhalf} in the article (source0.tex lines 852--868, source label \texttt{th:presymp-Hhalf}), delivered together with the article's own complete original proof of it (source0.tex lines 870--900, one \texttt{proof} environment of 31 lines). No mathematics is written, repaired or re-derived: statement and proof are the article's own text line for line. Required context retained uncounted: none. Every definition, display and label of those lines is retained unchanged. Omitted from this package and not redistributed: 1--851, 869--869, 901--1756. Nothing inside a retained excerpt is omitted or altered: every retained body is delivered exactly as the article's lines are written, so no omission receipt of any kind accompanies this package. Citations: the retained excerpts carry no citation command at all, so no bibliography entry of the article is transcribed and no reference is redistributed. Reference adaptation: none was needed and none was made; every reference inside the delivered unit resolves inside it, so no unresolved reference or citation is produced. Printed slips kept literal: no printed word, symbol, hypothesis or number of the counted statement or of its proof was altered. Layout and class: the article's own class is \texttt{amsart} with packages including amsmath, amsthm, mathrsfs, graphicx, hyperref, enumitem, mathabx, txfonts, array, tikz, algorithm2e; the wrapper is the lane's pinned \texttt{amsart} editorial wrapper at 10pt on A4, which restates the theorem-like environments the retained text needs and adds the article's own macro definitions that the retained text needs (none), with extra wrapper packages (bm, bbm, dsfont, xspace, cancel, mathtools). The delivered numbering is the unit's own. Assets: no retained span contains a figure environment, a graphics-inclusion command, a PDF-page inclusion, a table or a TikZ picture; no image file is copied into this package, the package declares no asset dependency and no image file is redistributed with it. Licence: version arXiv:2510.19830v1 of this article is distributed under the Creative Commons Attribution 4.0 International licence, as recorded in the arXiv licence metadata for this exact version and shown on the abstract page https://arxiv.org/abs/2510.19830v1. A full rights notice is delivered at licenses/arxiv-2510.19830-CC-BY-4.0.txt. Characters: every retained excerpt is pure ASCII, so the delivered engine, XeLaTeX with its default Latin Modern text font, reports no missing character.
\begin{theorem}[Canonical presymplectic form on $H^{1/2}$]\label{th:presymp-Hhalf}
On the Hilbert loop space $H^{1/2}(\mathbb S^1;\mathbb R^m)$,
the differential of the 1-form $\Theta_0$ is the continuous bilinear, skew form
\[
\boxed{\ d\Theta_0(\gamma)[h,k]\;=\;
\langle \partial_t h,\ k\rangle_{H^{-1/2},\,H^{1/2}}
-\langle \partial_t k,\ h\rangle_{H^{-1/2},\,H^{1/2}}\ ,\ }
\]
which may be written equivalently as
\[
d\Theta_0(\gamma)[h,k]\;=\;2\,\langle \partial_t h,\ k\rangle_{H^{-1/2},\,H^{1/2}}
\;=\;2c\int_{\mathbb S^1}\!\langle |D|^{1/2}h,\ \mathcal H\,|D|^{1/2}k\rangle\,dt,
\]
where $\mathcal H$ is the Hilbert transform and $|D|^{1/2}$ the fractional derivative.
In particular, $d\Theta_0$ is independent of $\gamma$ (translation-invariant) and defines
the \emph{canonical presymplectic form} on the loop space.
\end{theorem}

\begin{proof}
First assume $\gamma,h,k\in C^\infty(\mathbb S^1;\mathbb R^m)$.
Since the underlying space is linear, we may regard $h,k$ as constant vector fields.
Then
\[
\frac{d}{d\varepsilon}\Big|_{\varepsilon=0}\Theta_0(\gamma+\varepsilon h)[k]
=\int_{\mathbb S^1}\!\langle h'(t),k(t)\rangle\,dt,\qquad
\frac{d}{d\varepsilon}\Big|_{\varepsilon=0}\Theta_0(\gamma+\varepsilon k)[h]
=\int_{\mathbb S^1}\!\langle k'(t),h(t)\rangle\,dt.
\]
Thus
\[
d\Theta_0(\gamma)[h,k]
=\int\langle h',k\rangle-\int\langle k',h\rangle
=2\int\langle h',k\rangle
=-2\int\langle h,k'\rangle,
\]
where we used integration by parts on $\mathbb S^1$.
This shows skew-symmetry and independence of $\gamma$.

For general $h,k\in H^{1/2}$, approximate in $H^{1/2}$ by smooth sequences
$h_n,k_n$; use the boundedness of $\partial_t:H^{1/2}\to H^{-1/2}$ and continuity of the
duality to pass to the limit, which yields
\[
d\Theta_0(\gamma)[h,k]
=\langle \partial_t h,\ k\rangle_{H^{-1/2},H^{1/2}}
-\langle \partial_t k,\ h\rangle_{H^{-1/2},H^{1/2}}.
\]
Finally, the Fourier identity $\partial_t=\mathcal H\,|D|$ on mean-zero components
gives the last expression with $|D|^{1/2}$ and $\mathcal H$.
\end{proof}

\end{document}
